\documentclass[11pt]{article}
\usepackage[a4paper,margin=23mm]{geometry}
\usepackage{amsmath,amssymb,amsthm,booktabs,longtable}
\usepackage[colorlinks=true]{hyperref}
\newtheorem{proposition}{Proposition}
\newtheorem{lemma}{Lemma}
\newcommand{\ind}{\mathbf1}
\newcommand{\HH}{\mathcal H}
\newcommand{\hh}{\mathcal h}
\newcommand{\aa}{\mathfrak a}
\newcommand{\CN}{\mathcal C(N)}
\title{The own Chen weight and the actual 248-generation table}
\author{Independent insertion-chain audit}
\date{9 October 2026}
\begin{document}
\maketitle

This note directly checks the quotient weight, extends the old 24-generation
insertion proof to the actual 248-generation table, and derives the exact
term displaced when the lower starting cutoff moves. It does not infer
analytic admission from a numerical PASS. In particular,
\texttt{global-chen-weight-audit-stage.json} is an incomplete older audit;
it is not used as a proof. The later complete case proof in
\texttt{independent-global-chen-weight-audit.tex} was compared directly
with the present main-file definitions. We give its necessary argument
below. All old files, including the main manuscript, are unchanged.

The external standard linear-sieve, Bombieri--Vinogradov and dispersion
theorems are accepted as stated. Their finite-atom reductions, moving
residue argument and the positive product theorem are specified below.
We found two documentation gaps rather than a counterexample: the old
insertion theorem expressly stopped at 24 generations; and the low-rank
cutoff in the actual-$K$ audit is too restrictive to quote verbatim for
some one-prefix terms. Both gaps are repaired here. The ordinary causal
starting-tree compiler is a separate argument; this note neither assumes
that a numerical density is automatically such a tree nor reaudits that
compiler. Its newly completed, separately bound statement is
\texttt{own-lower-compiler-chain-audit-20261009.tex}; it supplies the
ordinary causal start at $w=a,h$, not a replacement for the table-child
gate checks proved here.

\section*{1. The actual pointwise inequality}
Fix
\[
 a=417/5000,\quad h=19/125,\quad d=59/250,
 \quad 0<a<h<d<1/3,\quad h+d>1/3,\quad4a+d<1.
\]
For $1\le n\le N$ put
$Q_b(y;n)=\ind_{b\mid n}\ind_{P^-(n/b)\ge y}$, with $P^-(1)=\infty$.
Prefix primes in every sum are distinct. The eleven functions are
\begin{align*}
Y_1&=Q_1(N^a),&Y_2&=Q_1(N^h),\\
Y_3&=\sum_{N^a\le p<N^{1/3}}Q_p(N^a),&
Y_4&=\sum_{N^a\le p<N^d}Q_p(N^a),\\
Y_5&=\sum_{N^a\le p_1<p_2<N^h}Q_{p_1p_2}(N^a),&
Y_6&=\sum_{N^a\le p_1<N^h\le p_2<N^d}Q_{p_1p_2}(N^a),\\
Y_7&=\sum_{N^d\le p_1<p_2<(N/p_1)^{1/2}}Q_{p_1p_2}(p_2),\\
Y_8&=\sum_{N^a\le p_1<N^{1/3}\le p_2<(N/p_1)^{1/2}}Q_{p_1p_2}(p_2),\\
Y_9&=\sum_{N^h\le p_1<N^d\le p_2<(N/p_1)^{1/2}}
 Q_{p_1p_2}((N/p_1p_2)^{1/2}),\\
Y_{10}&=\sum_{N^a\le p_1<p_2<p_3<p_4<N^h}Q_{p_1p_2p_3p_4}(p_2),\\
Y_{11}&=\sum_{N^a\le p_1<p_2<p_3<N^h\le p_4<N^d}Q_{p_1p_2p_3p_4}(p_2).
\end{align*}
There is no extra restriction on $p_3p_4$ in $Y_{11}$.
Let $W=3Y_1+Y_2-Y_3-Y_4+Y_5+Y_6-2Y_7-Y_8-Y_9-Y_{10}-Y_{11}$.

If $n$ is not $N^a$-rough all terms vanish. For $Y_9$ this follows from
$p_1p_2<N^{(1+d)/2}$ and its quotient cutoff
$>N^{(1-d)/4}>N^a$; for every other term it is immediate.
For squarefree rough $n$ write $q_1<\cdots<q_k$ and let $r,s,t$ count
the factors below $N^h,N^d,N^{1/3}$, respectively. Then
\[
(Y_1,\ldots,Y_6)=(1,\ind_{r=0},t,s,{r\choose2},r(s-r)).
\]
A quotient tested at its second selected factor forces the first two
selected factors to be $q_1,q_2$. Thus, for $r\ge2$,
$Y_{10}={r-2\choose2}$ and $Y_{11}=(r-2)(s-r)$, and consequently
$W\le s-t\le0$. The convention ${0\choose2}={1\choose2}=0$ includes
$r=2,3$. For $k\ge3$, $q_1q_2^2<q_1q_2q_3\le N$; hence
$Y_7=\ind_{s=0}$, $Y_8=\ind_{t=1}$, and $t\ge1$.
The remaining cases give the following upper bounds:
\[
\begin{array}{c|c}
r=1&2-t-\ind_{t=1}\\
r=0,\ s\ge2&4-s-t\\
r=s=0&2-t-\ind_{t=1}\\
r=0,\ s=1,\ t\ge3&3-t\\
r=0,\ s=1,\ t\le2&3-t-1-\ind_{t=1}.
\end{array}
\]
The last row uses $Y_9\ge1$, not $Y_9=1$: if $k\ge4$, then
$n\ge N^{h+d+2/3}>N$, so $k=3$, and
$q_1q_2q_3^2\ge N^{h+d+2/3}>N$ proves that $(q_1,q_2)$ is one
admissible $Y_9$ pair. Every displayed bound is nonpositive.
If $k\le2$, either $Y_5+Y_6=0$ and $W\le4$, or this sum is one,
$Y_2=0$, $Y_3=Y_4=2$, and $W\le0$. The unit has $W=4$.
This is an exhaustive analytic proof of
\[
                  W(n)\le4\ind_{\Omega(n)\le2}
                  \quad(n\text{ squarefree},\ 1\le n\le N).
\]
No assumption $n\ge N^{1-a}$ was used. Actual $Y_9$ multiplicities
are retained in its subsequent upper estimate.

All nonzero terms have at most $M=\lfloor1/a\rfloor=11$ prime factors.
For rough nonsquarefree integers the elementary bound
\[
 |W|\le4+2M+6{M\choose2}+2{M\choose4}=1016
\]
suffices. The number of integers $n\le N$ divisible by $p^2$ for
some $p\ge N^a$ is at most $2N^{1-a}$ for sufficiently large $N$.
Consequently, on summing over primes $P<N$ with $(P,N)=1$,
\begin{equation}\label{eq:eleven}
4D_{1,2}(N)\ge3\Upsilon_1+\Upsilon_2+\Upsilon_5+\Upsilon_6
 -\Upsilon_3-\Upsilon_4-2\Upsilon_7-\Upsilon_8-\Upsilon_9
 -\Upsilon_{10}-\Upsilon_{11}-2032N^{1-a}.
\end{equation}
Here $\Upsilon_j=\sum_PY_j(N-P)$, and this exact inequality uses
$\Omega(1)=0$, as in the main manuscript. If $D_{1,2}$ is instead
defined with $N-P>1$, subtract a further $4$ on the right of
\eqref{eq:eleven}; there is at most one excluded prime $P=N-1$.
This is $O(1)$ and does not change any normalized conclusion.
The error is
$o(\CN N/\log^2N)$ uniformly in even $N$; $\CN$ is bounded below
by the fixed positive twin-prime product. Since sufficiently large even
$N$ is not prime, $N-P=0$ never occurs. Restricting to $(P,N)=1$
only reduces the left-hand representation count.

\section*{2. What a table estimate means on a real prime family}
Put
\[
 A(s)=sF(s)/(2e^\gamma),\qquad \aa(s)=sf(s)/(2e^\gamma),\qquad
 A=1\ (1\le s\le3),\quad\aa=0\ (s\le2).
\]
Fix a nonnegative finite marked prime-box weight $\sigma(b)$, with
$b=N^{T+o(1)}$, one total level $Q=N^\nu$ and residual exponent
$L=\nu-T>0$. The same-level normalization is
\[
 \Theta_Q=4\operatorname{li}(N)
       \sum_b{\sigma(b)\mathcal C(bN)\over\varphi(b)\log(Q/b)}.
\]
An upper defect $H(s)$ means an upper count $(A(s)-H(s)+o(1))\Theta_Q$;
a lower defect $h(s)$ means $(\aa(s)+h(s)+o(1))\Theta_Q$.
These are estimates averaged over the indicated prime boxes, not
estimates for each fixed integer $b$.

At every enhanced state retain the marked slots, a branch, its strict
cap margin, all full square tests
\[
 \sum_{i<j}t_i+2t_j<\nu,
\]
all selected-remaining tests, and a physical ceiling $\chi\le\min t_i$.
A future-third-slot bound $c$ also requires $\chi\le c$. The actual
comparison cutoff must be at most $N^\chi$, and
$\alpha=(1-T)/L\ge2$. In a classical state choose $\nu<1/2$ with
margin and use the absolute complete-modulus BV bound instead.

The ordinary residual modified weights are constructed at $Q/b$;
they are not restrictions of a global weight at $Q$. Their finite
rough-box expansion retains its signed scalar, small-prime weight
$\beta$ and original repeated-label factorial. Appending the prefix
concatenates its actual full-square tests with the residual well
support. For a selected subset retain its original slots; the additional
tests, after the old slots have occurred, are the same full tests.
The selected-subproduct factorization lemma is applied only with both
sets of tests, not just the size of the selected product.

For each buffered atom, the endpoint version of L23 Proposition 6.1
and its proved selected-product extension supply a fixed partition of
the labelled rough slots. Dividing each group convolution by its
group factorial makes it one-bounded; multiplying the original scalar
by those same factorials is an exact identity. The original
repeated-label factorial is not inserted twice. Appending $\beta$
preserves the bound by uniqueness of the small/rough decomposition.
L23 Proposition 5.2 then applies to each fixed atom with either sign.
This is why signed feedback is licensed; arbitrary functions supported
on well moduli are not being called well-factorable.

\begin{lemma}[Low ranks for the actual table families]\label{lem:lowrank}
All enhanced table families in Section 4, and their finite descendants,
have largest actual rough-prime exponent $x<e=57/224$ and total level
$\nu\le\Lambda=153/256$. An atom with at most two rough primes,
including its small factor, has support below $N^{1/2-c}$ for fixed
positive $c$ after choosing the construction precision sufficiently
small on a fixed buffered box.
\end{lemma}
\begin{proof}
For rank two write $x\ge y$. Full well support gives $x+2y\le\nu$
up to a precision loss; therefore
\[
 x+y\le(\nu+x)/2\le(\Lambda+e)/2
 =1527/3584<1/2,\qquad
 1/2-1527/3584=265/3584.
\]
Ranks zero and one are easier. The $N^{O(\eta)}$ small factor and
interval enlargement are absorbed in this fixed gap. All descendants
append primes below the previous smallest slot, so the largest slot
never increases. Classical states already have complete support below
half level. Boundedly many labelled prefix representations and the
fixed construction coefficients give a bounded complete-modulus
coefficient, to which absolute BV applies.
\end{proof}
The $19/96$ low-rank cutoff printed in the actual-$K$ distribution
audit is specific to that strategy. It is not needed, and is not
silently assumed, in this lemma. For rank at least three the retained
largest, selected-product or $W$ branch supplies the factorization.

The moving residue is $a_{\rm res}=N$. On prime-variable scales
$x\ge N^{1-\kappa}/O(1)$, take $\kappa$ after the fixed factorization
margin so that $N<x^{1+\epsilon}$, as required by L23 Proposition 5.2.
The remaining scale perturbations are $O(\kappa)$. The omitted small
prime-variable tail is at most $N^{1-\kappa+o(1)}$ by the divisor
bound. For a rough collision $p^2\mid d$, $p>N^\omega$, the full
HB-weighted progression and exact-coprime mean error is
$\ll(N+N^{\theta_*})N^{o(1)-\omega}$, $\theta_*<3/5$.
Thus no unweighted collision estimate is used after HB expansion.
These are the explicit signed-atom and moving-residue reductions in
the main prefix interface and the later joint distribution repair.
The presieve retains the full $N$-dependent local product, including
the omitted primes dividing $N$.

Uniformly on fixed positive boxes,
$\mathcal C(bN)/\mathcal C(N)=1+o(1)$ and
$\varphi(b)/b=1+o(1)$, and deletion of prefix primes dividing $N$
has $o(1)$ reciprocal mass. Hence dividing by
$\mathcal H_N=\CN N/\log^2N$ converts $\Theta_Q$ into
\begin{equation}\label{eq:norm}
             4\int {d\mu(\boldsymbol t)\over L},
 \qquad d\mu=\prod_i{dt_i\over t_i}.
\end{equation}
This fixes every factor four below. Appending $N^{Lu}$ gives the
normalized child mass $du/[u(1-u)]$; appending $N^{Lu},N^{Lv}$ gives
$du\,dv/[uv(1-u-v)]$. With $w=(1-u)/u$, the first becomes $-dw/w$.

\section*{3. Finite inheritance through the actual 248 generations}
The lineage is $24\to56\to120\to184\to248$, respectively adding
$32,64,64,64$ finite rounds. Every table has 200 closed cells
$[1+i/50,1+(i+1)/50]$, denominator $2^{80}$. Values are step lower
bounds, not interpolants. A fixed point can use either containing
closed cell; a whole box uses a bound valid over its entire argument
range, or zero outside support. Endpoint comparisons are obtained by
buffering the level and then taking its limit. At a clamped argument
two the cutoff is structurally $L/2$ and is handled as such, not
discarded as a zero-measure set.

The initial seed is a proved finite four-round W2 construction with
$K=6.81$, uniform for $\alpha\ge2$ on the admitted states. It gives
$H=2327/100000$ on $[1,11/5]$ by positive-count cutoff comparison.
Below two this does not require a new ordinary theorem: comparison
to argument $11/5$ lowers the physical cutoff. The same device, with
the full $A(s)-A(v)$ term where needed, licenses later upper comparison.

\paragraph{Buchstab.}
For $s\ge2$ the finite actions use
\[
 H_{\rm B}(s)=\int_{\max(2,s-1)}^5h(w){dw\over w},\qquad
 h_{\rm B}(s)=\int_{s-1}^5H(w){dw\over w}.
\]
They start with the ordinary baseline at argument six. The child prime
has exponent $x=L/(w+1)$, residual $Lw/(w+1)$ and
$\alpha'=\alpha+(\alpha-1)/w\ge2$. Since $w\ge1$, $2x\le L$;
the equality boundary is approached from strict interiors. Also
$x\le L/s\le\chi$ and the new ceiling is $x$. In the lower formula
the range $1<w<2$ is retained: its upper estimate uses smaller-cutoff
comparison and $A(w)=1$. Truncating that integral at two would lose a
valid positive contribution. Closed target cells use their right endpoint
in the lower integration limit, giving a uniform, conservative bound.

\paragraph{The 70 W1 actions.}
They have target $12/5\le s\le3$ and $16/5\le s_0\le5$.
Their $H(s_0)$ term retains the prefix. The lower child has
$u\in[1/s_0,1/s]$, residual $L(1-u)$, physical cutoff $L/s_0$,
argument $s_0(1-u)$, and
$\alpha'=(\alpha-u)/(1-u)\ge2$. The new prime $Lu$ is below
$L/s\le\chi$, satisfies $2u<1$, and exceeds the physical child
cutoff. Arguments below two contribute only nonnegativity. The switched
triple is handled by the separate positive product theorem, with the
uniform rough-mass bound and coefficient $K/8$; it is not a signed
ordinary child. The final bank retains $K=6.7633$ for W1.

\paragraph{All 36 W2 actions, including the 33 additions.}
Write their exact parameters as $(s,s_0,\phi_1,\phi_2,\phi_3)$ and
$a_0=1/s_0$, $a_i=1/\phi_i$, $b=1/s$. The companion check independently
checks every rational tuple, without importing the author validator:
\[
2<s<3,\quad s<\phi_3<\phi_2<\phi_1<s_0\le5,
\quad a_3+2a_1<1,\quad3a_1<1,\quad a_2+2a_1<1,
\]
\[
a_0<(1-a_2)/2<1/3<a_2<1/2,
\quad s_0(1-b)>2,\quad s_0-1\le4,
\quad s_0(2-2a_2)>2.
\]
The unchanged-prefix terms are $4H(s_0)+H(\phi_1)$. The three
negative singleton terms have $u\in(a_0,b),(a_0,a_2),(a_0,a_3)$,
child argument $s_0(1-u)\in(2,4]$, physical cutoff $La_0$, and the
same singleton square and ratio calculation as W1.

For two new primes $Lu<Lv$ the domains and actual cutoffs are
\[
\begin{array}{c|c|c}
\Gamma_5^+&a_0<u<v<a_2,\ v+2u<1&La_0\\
\Gamma_6&(a_0,a_1)\times(a_2,a_3)&La_0\\
\Gamma_7&a_0<u<v<a_1&Lu\\
\Gamma_8&(a_0,a_1)\times(a_1,a_2)&Lu.
\end{array}
\]
Both new full-square tests are $2v<1$ and $v+2u<1$.
The displayed parameter tests prove these on all four retained domains.
Old selected-remaining tests are unchanged; the new ones are exactly
these full tests after multiplying by the old prefix. The actual child
ratio is $(\alpha-u-v)/(1-u-v)\ge2$. The arguments are
$s_0(1-u-v)$ for terms 5,6 and $(1-u-v)/u$ for terms 7,8.
They lie in $(1,s_0-2]\subset(1,3]$. This follows from
$1-u-v>u\ge a_0$; their maxima occur at the lower endpoints.
For arguments below two, comparison lowers the child cutoff; it never
raises it above $Lu$, the child's smallest specified prime.

The bad part of $\Gamma_5$ has no ordinary $H$ feedback. Its old
ordinary mass is removed once and replaced by a positive product upper
bound. Terms $\Gamma_9$--$\Gamma_{21}$ likewise use positive product
estimates, with the uniform envelope $\omega(x)\le\max(1/x,0.567144)$
for $x\ge1$, and its conservative extension equal to one below one.
Neither the nonunit domains nor the unit surfaces in terms 20,21 are
dropped by setting the parent ratio equal to two. The current W2 cost
uses $K=6.7619$. Its rigorous variable-kernel saving subtracts a smaller
upper cost from that same cost; it is not extra child feedback.

Thus each W2 candidate is exactly a certified ordinary/switching base
plus
\[
 {1\over5}\left\{4H(s_0)+H(\phi_1)
 +\sum_{\Gamma_2,\Gamma_3,\Gamma_4}\int h(s_0(1-u)){du\over u(1-u)}
 +\sum_{i=5^+,6,7,8}\iint H(\rho_i){du\,dv\over uv(1-u-v)}\right\}.
\]
The division by five occurs once. All feedback measures are positive.
The $A(s)=1$ baseline is valid because all action targets have $s<3$.
An action transfers only to closed cells with right endpoint at most
its target $s$; positive-count comparison then lowers the cutoff.

\paragraph{Why this is an induction on real families.}
At every ordinary child all new primes are below the old ceiling,
the full and selected square tests just proved hold, and any retained
future-third ceiling persists. The total level stays $Q$; the actual
$\alpha$ increases and is never reset to two. Therefore Section 2
applies to that child. At switched children only the positive product
theorem is invoked. That theorem allows every fixed positive factor
cutoff, by choosing its HB precision after the finite family is fixed.
The minimum of $1-2/s$ over the 36 actions is positive, and BH leaves
at least half the residual exponent. Each appended prime has exponent
at least $L/6$ relative to its parent. Any fixed finite recursion tree
therefore has a positive minimum exponent. There is no assertion
uniform in recursion depth or in a cutoff tending to zero.

Every round uses only the preceding round's bounds. Taking a finite
maximum of valid estimates at the same state, rounding positive weights
down, and taking minima on cell-crossing boxes preserve validity.
This proves the 248-generation inheritance. The numerical integrations,
all 33 new ordinary collections, all 5224 pair weights and all 64 new
400-integer arrays are separately covered by
\texttt{expanded248-numeric-closure-independent.json}; those numerical
facts are used only after the foregoing actual-state proof.

\section*{4. The eleven occurrences and their coefficients}
Use $\mu=25/128$, $e=57/224$, $\Lambda=153/256$,
$\nu_*=249/448$, $w_*=25/224$ and
$v(t)=\min((1+t)/2,57/64-3t/2)$. Let $\HH,\hh$ be exactly the
248-generation step table. The original coefficients of
$(L_1,L_2,L_5,L_6,U_3,U_4,U_7,U_8,U_9,U_{10},U_{11})$
are $(3,1,1,1,-1,-1,-2,-1,-1,-1,-1)/4$.

\paragraph{$L_1$.}
Its ordinary value is $4\aa(\nu_*/a)/\nu_*$. There is no direct
$\hh$ insertion into this empty-prefix state: $1/\nu_*<2$.
Any new ordinary causal starting-tree improvement has coefficient
$3/4$ and is separately compiled. Nothing from the 248 table is
subtracted when that start is improved.

\paragraph{$L_2$.}
In the subtracted one-prefix family $w_*<t<h$, use
$\nu=(1+t)/2$, $L=(1-t)/2$, $\chi=t$ and
$s=(1-t)/(2t)\in[53/19,199/50]$. Then $2t<\nu$, $\alpha=2$
and the physical cutoff is exactly $t$. Later sieve primes are strictly
smaller; a structural equality is allowed. The $H$ term in the
unquartered $L_2$ is $4\int_{w_*}^h\HH(s)\,dt/(tL)$.
It contributes positively to the Chen coefficient.

\paragraph{$U_3,U_4$.}
Their domains are $(a,1/3)$ and $(a,d)$. Use $\nu=v(t)$ for
$t<e$; the $U_3$ tail uses the limit $\nu\uparrow1/2$ from below.
Put $L=\nu-t$ and compare the original cutoff $a$ down to
$\min(a,L/2)$, giving $s=\max(2,L/a)$. The enhanced branch has
$2t<v(t)$, $\alpha\ge2$ and ceiling $t$. The classical tail needs
no enhanced square condition. Thus the normalized upper bound is
$4\int(A(s)-\HH(s))dt/(tL)$; $\HH=0$ outside $[1,5]$.
Both occurrences have coefficient $-1/4$, so both defects improve
the final lower bound. The clamp does not change the ordinary
multiplier because $A=1$ below three.

\paragraph{$L_5,L_6$.}
Use decreasing slots $t>u$, domains $a<u<t<h$ and $a<u<h<t<d$,
and
\[
 V=\begin{cases}\max(v(t),v(t+u),\Lambda-3a/8),&t<\mu,\\
 v(t),&t\ge\mu,\end{cases}\qquad L=V-t-u,\quad s=L/a.
\]
On $s\ge2$, $2t<v(t)\le V$ and $t+2u<V$ since $2a>h\ge u$.
The largest and leading-pair branches retain ceiling $u$; the constant
$W$ branch retains ceiling $a$ and its future-third bound $a$.
It is used above its $v(t)$ gate only if it improves that branch.
The ratio is at least two: $2v(t+u)\le1+t+u$, and for the $W$
branch $a\ge25/352$ gives $\Lambda-3a/8\le1/2+a\le(1+t+u)/2$.
The argument is at most $1/(2a)-1=2083/417<5$.
The bound is
$4\iint(\aa(s)+\hh(s))\,dt\,du/(tuL)$ on the admitted region;
below two use zero. Each occurrence has coefficient $+1/4$.

\paragraph{$U_7,U_8,U_9$.}
No table defect occurs here. Their exact marked positive-product sums
are bounded with $K=6.7619$. The moving square-root cutoff and every
$Y_9$ label remain present. In particular replacing $Y_9$ by an indicator
would not justify the original negative term. Their coefficients are
$-1/2,-1/4,-1/4$, respectively. These are separate from positive
product costs already used to prove the table; reusing a valid theorem
in two estimates does not spend its constant twice.

\paragraph{$U_{10},U_{11}$.}
Use decreasing slots $x>y>z>w$ with domains
$a<w<z<y<x<h$ and $a<w<z<y<h<x<d$. The actual quotient cutoff
is $z$. The ordinary competing estimate compares down to
$\min(w,L/2)$, with $L=\nu-x-y-z-w$.
Set $W_4=\max(2x,x+2y,x+y+2z,x+y+z+2w)$.
For a selected subset $I$ require
\[
 v(\sum_{i\in I}t_i)>\max(1/2,W_4,G_I),\quad
 G_I=\max_{j\notin I}\left(\sum_{i\in I}t_i+
                  \sum_{k<j,\,k\notin I}t_k+2t_j\right).
\]
The extra $W$ cap $\Lambda-3z/8$ requires $x<\mu$ and that it
exceed $1/2,W_4,v(x)$. Freeze one licensed branch and a strictly lower
level; retain ceiling $w$. Its first three slots cannot change.
Also $\alpha\ge2$ because $x+y+z+w>4a>2\Lambda-1$.
For a selected subset of size one or two, $v(S_I)\le(1+S_I)/2$
gives $L/w\le1/(2a)-7/2$ or $1/(2a)-3$. For size three or four,
$S_I\ge3a>\mu$ puts it on the descending branch and gives
$L/w\le57/(64a)-17/2$ or $57/(64a)-10$. The $W$ cap gives
$L/w\le\Lambda/a-35/8$. The maximum of these five bounds is
$1249/417<3$, so $s=\max(2,L/w)<3$.
For each common positive box the upper option is
\[
 \min\left\{K{\omega((1-x-y-z-w)/z)\over z},
                  {4(1-\HH(s))\over L}\right\}d\mu_4.
\]
Both occurrences have coefficient $-1/4$. The two options are a
minimum on the same positive family, never additive savings on
overlapping covers. Repeated-bin factorials belong to the original
ordered-domain measure and are retained once.

\section*{5. Replacing the L2 start removes exactly the same table term}
For any fixed $w\in[w_*,h]$, the least-prime identity is
\[
 S(\mathcal A,N^h)=S(\mathcal A,N^w)
                  -\sum_{N^w\le p<N^h}S(\mathcal A_p,p).
\]
It is an identity of positive integer counts, not an identity obtained
by subtracting sieve asymptotics with unrelated levels. Let $I(w)$ be
the separately proved extra normalized density of the ordinary
starting lower tree at cutoff $w$ and the original level $\nu_*$.
Define
\[
 B(w)={4\over\nu_*}\aa(\nu_*/w)
          -8\aa((1+w)/(2w)).
\]
The ordinary upper kernel is $8A((1-t)/(2t))/[t(1-t)]$.
The delay identity $\aa'(q)=A(q-1)/(q-1)$ gives
\[
 {d\over dt}\aa((1+t)/(2t))
       =-{A((1-t)/(2t))\over t(1-t)}.
\]
It follows without replacing $F,f$ by truncated explicit formulas that
the new normalized $L_2$ is
\[
 B(w)+8\aa((1+h)/(2h))+I(w)
        +4\int_w^h {\HH((1-t)/(2t))\over t(1-t)/2}\,dt.
\]
At $w_*$, $\nu_*=(1+w_*)/2$, giving precisely the old ordinary
$8\aa((1+h)/(2h))-4\aa(249/50)/(249/50)$ term. Therefore the
replacement increment in the \emph{complete Chen coefficient} is
\begin{equation}\label{eq:move}
 {B(w)-B(w_*)+I(w)\over4}
       -\int_{w_*}^w {2\HH((1-t)/(2t))\over t(1-t)}\,dt.
\end{equation}
At $w=h$ the entire one-prefix $H$ integral disappears. It must be
the 248 table's integral, with its same cell convention and same
normalization. It is neither the 184 integral nor the initial seed
integral. Adding $I(w)/4$ alone would be incorrect. Other table
occurrences remain unchanged. An $L_1$ starting-tree improvement has
coefficient $3/4$ and has no displaced $H$ term.

The saved baseline248 is a lower bound for the explicit eleven-term
functional with these term choices, rather than a black-box count
constant unrelated to them. Hence an independently bounded difference
of these explicit terms can be added to its saved lower bound.
Using an upper bound on the removed positive $H$ integral is the
conservative direction. This justifies the substitution mechanism;
the actual density of the new causal starting tree is established
separately, not by this algebra.

\section*{6. Quantifiers and audit boundary}
First fix the finite depth, positive exponent cutoffs, branch choices,
closed-cell prescriptions and desired error. Choose a common finite
marked partition and strict factorization buffers, then a sufficiently
fine residual sieve construction, and finally let $N\to\infty$.
Only after this take endpoint/exhaustion limits. Signed scalar sums
and group factorials are finite and independent of $N$; the external
remainder theorem is uniform in the moving residue. Boundaries can be
trimmed using positive majorants on fixed compact domains. Structural
cutoff equalities are handled by lower cutoff comparison and continuous
level limits, not discarded. No individual-prime distribution assertion,
infinite fixed point, uniform-in-depth error, or interpolation is used.

The local checks accompanying this note concern exact parameter gates,
lineage, point targets, normalization identities and bindings. They do
not rerun old large integrations. Existing full independent numerical
records are bound separately. No unresolved table-admission condition
was found for these actual parameters and actions. This conclusion
accepts L21's ordinary sieve construction, L23's stated distribution
and factorization theorems, the explicitly supplied selected-product
repair, and the established positive-product $K$ theorem. It is not a
new verification of their underlying spectral proofs, nor a claim about
the separate literature-priority objective or an effective threshold
for $N$.
\end{document}
